\documentclass{standalone}
\usepackage{pgfplots}
\pgfplotsset{compat=newest}
\usepgfplotslibrary{colormaps,patchplots}
\usepackage{amsmath}
\usepackage{tikz}
\usetikzlibrary{graphs, graphdrawing, quotes}
\usegdlibrary{circular}

\begin{document}
		
		\begin{figure}[h!]
			\centering
			\begin{tikzpicture}
				\graph[circular placement, radius=3cm, nodes={draw, circle, minimum size=1cm}] {
					1/"1 (3 min)",
					2/"2 (4 min)",
					3/"3 (5 min)",
					4/"4 (6 min)",
					5/"5 (7 min)",
					6/"6 (10 min)",
					
					% Aristas de las parejas óptimas
					1 -- 6,
					2 -- 5,
					3 -- 4
				};
			\end{tikzpicture}
			\caption{Grafo de las parejas óptimas para cruzar el río}
		\end{figure}
		
\end{document}
	
	
	\begin{tikzpicture}
		\begin{axis}[
			view={45}{30},
			axis lines=middle,
			xlabel={$x$},
			ylabel={$y$},
			zlabel={$z$},
			xtick={-2, -1, 0, 1, 2},
			ytick={-2, -1, 0, 1, 2},
			ztick={0, 1, 2, 3, 4, 5, 6},
			xmin=-2.5, xmax=2.5,
			ymin=-2.5, ymax=2.5,
			zmin=-1, zmax=6,
			domain=-1:2*pi,
			samples=40,
			samples y=20,
			colormap name=viridis,
			enlargelimits=false,
			axis line style={thick},
			tick style={gray},
			label style={font=\small},
			tick label style={font=\tiny, inner sep=0pt, yshift=-2pt}, % ← Ajuste de cercanía
			]
			% Superficie del cono con transparencia
			\addplot3[
			surf,
			shader=interp,
			opacity=0.4,
			domain=0:2*pi,
			y domain=0:5,
			]
			({2*(1-y/5)*cos(deg(x))}, {2*(1-y/5)*sin(deg(x))}, {y});
			
			% Secciones transversales
			\foreach \h in {0,1,2,3,4,4.5} {
				\addplot3[
				thin,
				color=gray,
				domain=0:2*pi,
				samples=100,
				]
				({2*(1-\h/5)*cos(deg(x))}, {2*(1-\h/5)*sin(deg(x))}, {\h});
			}
			
			% Fórmula del cono como anotación
			\node at (axis cs:0,0,6.2) [anchor=south] {$x^2 + y^2 = \left(1 - \frac{z}{5}\right)^2 \cdot 4$};
		\end{axis}
	\end{tikzpicture}
\end{document}